\chapter{Conclusions}
\label{ch:conclusion}

This dissertation asks what environmental and terrain factors explain why some Sitka spruce plots in Aberfoyle grow faster or slower than the Chapman-Richards baseline predicts. The approach: compute CR residuals across a cleaned, balanced cohort of plots surveyed by airborne LiDAR at up to six timestamps between 2002 and 2023; test whether the residuals cluster spatially; identify which terrain and wind features best explain them using XGBoost and SHAP; then build a physics-informed neural network in which the growth ceiling $y_{\max}$ is learned from those same features instead of fixed globally.

At this draft stage, the cleaned cohorts, the leakage-aware variable set, and the modelling pipeline are defined (Chapter~\ref{ch:data}). Terrain and wind feature extraction, the CR fit on the current cohorts, and all downstream modelling are outstanding (Chapter~\ref{ch:results}). The dissertation will deliver: a spatial residual analysis with Moran's I and trajectory classification; an XGBoost and SHAP attribution comparing terrain-only and terrain-plus-wind feature sets; an Env-PINN with terrain-conditioned and fully-conditioned versions, evaluated against a baseline PINN and simpler baselines; an ablation study isolating what drives any improvement; and a convergent validity check between the Env-PINN's learned feature sensitivity and the SHAP ranking.

The intended contribution is a spatial map of learned growth ceilings across Aberfoyle, showing which parts of the forest are environmentally constrained and by how much. This is framed as attribution evidence, not as a causal or predictive claim beyond what the data and validation can support.
